\documentclass{article}
\usepackage{hyperref}
\title{Preserve an old operation by checking which site it can affect}
\begin{document}
\maketitle
In this follow-up, canceling an old queued operation is sufficient
for rescue, but gives up useful old work. If the old operation
can affect only one verified site through the admission cut,
the same rescue and unrelated delivery succeed while the old
operation also makes native progress. A predicted site that can
change after the check is insufficient.

\paragraph{Native source of the scope.}
The warehouse controller reads the intended old forward target
from the pinned RWARE agent's request location, not from the
correct answer. The occupied berth produces actual blocking when
the old call runs. A legacy carrier stages an old shelf but still
owes its return. In the radio plant the controller reads the
sender of an actual queued SubT datagram. Its execution consumes
that sender's rolling allowance and delivers the old payload to
a distinct transit cache via a native callback. Final delivery of
that old payload is still owed and it remains in the original
buffer. Such staging is useful progress in this model, not
acceptance or withdrawal of the older obligation. The concurrent
rescue uses one dispatch grant; the unrelated job has its own.

\paragraph{Bound at commitment.}
The policy chooses a site outside the union of old effects that
are still permitted before rescue commitment. It uses the actual
scoped observation when the owned submission gate freezes the
operation, and both possible sites when rebind remains allowed
with no fresh observation. Each scope is relative to the named
rescue resource, not the coverage of all hardware effects.
A passed check needs authority to hold that interpretation
through commit. Neither source supplies a general atomic transfer
API; the fixture gives the gate responsibility for its own queue.

The rebind control changes the warehouse carrier's heading with
native turns without moving it, or replaces the pending datagram
before dispatch. A stale scope directs rescue to the new blocked
site. Rejecting the now-mismatched old action at the gate recovers
the rescue but loses old progress. Checking the new scope at
commit and taking the other site recovers both; this case
assumes a timely observation and no intervening authority change.

\paragraph{Executed contrast.}
Both backends run 36 cases. With the scope held, both due
deliveries and the old staging effect occur. Under blanket
retirement both due deliveries occur but no old staging does.
With unobserved expressible rebinding, preserving the old grant
makes any single fixed rescue lose a branch; the other job and
old staging can still proceed. When the committed scope is known,
the matched direct policy attains the same result. This improves
on blanket retirement in the controlled workload at identical
grant count; it is not an improvement over a correctly scoped
direct architecture, a physical deadline, or an independent
I1 or I2 result.
\end{document}
